\chapter{应力状态分析和强度理论}

\intro{你的眼睛还没掉转过来看我，只起了一个势，我早惊乱得同一只听到弹弓弦子响中的小雀了。
\rightline{——沈从文}}

\section{研究应力状态的方法——单元体法}

\subsection{单元体的概念}
为了研究构件内某一点处的应力状态，我们围绕该点截取一个无限小的正六面体，称为\textbf{单元体 (Elementary Block)}。
\begin{itemize}
    \item \textbf{微元性}：各面上的应力被认为是均匀分布的。
    \item \textbf{平行性}：相对的两个面上的应力大小相等、方向相反（满足平衡条件）。
\end{itemize}

\subsection{应力的符号规定}
\begin{enumerate}
    \item \textbf{正应力 $\sigma$}：拉应力为正（离开截面），压应力为负（指向截面）。
    \item \textbf{切应力 $\tau$}：使单元体或其局部产生顺时针转动趋势的切应力为正，反之为负。
\end{enumerate}

\begin{myTheorem}{切应力互等定理}
    在相互垂直的两个平面上，切应力必然成对出现，且大小相等，方向垂直于两平面的交线，同时指向或同时背离该交线。
    \[
    \tau_{xy} = \tau_{yx}
    \]
\end{myTheorem}

\section{解析法求解二向应力状态}

\subsection{任意斜截面上的应力推导}

\begin{figure}[H]
    \centering
    \begin{tikzpicture}[scale=2.8, font=\small]
        \def\ang{30} 
        \coordinate (O) at (0,0);
        \coordinate (P) at (0,{1.5*tan(\ang)}); 
        \coordinate (Q) at (1.5,0); 
        
        \draw[thick, fill=structurecolor!5] (O) -- (Q) -- (P) -- cycle;
        
        \coordinate (M) at (0.75,{0.75*tan(\ang)}); 
        \draw[dashed, gray] (M) -- +(0.8,0) node[right]{$x$};
        
        \draw[-{Stealth}, color=second, ultra thick] (M) -- +(\ang:0.7) node[above right] {$\sigma_\alpha$};
        \draw[-{Stealth}, color=third, ultra thick] (M) -- +({\ang-90}:0.5) node[below right] {$\tau_\alpha$};
        
        \draw[->] ([shift=(0:0.4)]M) arc (0:\ang:0.4);
        \node at ($(M)+(0.5,0.15)$) {$\alpha$};
        
        \draw[-{Stealth}, color=main] (-0.4, {0.75*tan(\ang)}) -- (0, {0.75*tan(\ang)}) node[left] {$\sigma_x dA \sin\alpha$};
        \draw[-{Stealth}, color=main] (0.75, -0.4) -- (0.75, 0) node[below] {$\sigma_y dA \cos\alpha$};
    \end{tikzpicture}
    \caption{二向应力状态隔离体受力分析}
\end{figure}

设斜截面面积为 $dA$，由隔离体的静力平衡条件 $\sum F_n = 0$ 和 $\sum F_t = 0$ 可得：

\begin{enumerate}
    \item \textbf{法向平衡方程：}
    \begin{equation*}
        \sigma_\alpha dA = \sigma_x (dA \sin\alpha) \sin\alpha + \sigma_y (dA \cos\alpha) \cos\alpha + \tau_{xy} (dA \sin\alpha) \cos\alpha + \tau_{yx} (dA \cos\alpha) \sin\alpha
    \end{equation*}
    利用切应力互等定理 $\tau_{xy} = \tau_{yx}$，并引入倍角公式，化简为：
    \begin{myformula}
        \sigma_\alpha = \frac{\sigma_x + \sigma_y}{2} + \frac{\sigma_x - \sigma_y}{2} \cos 2\alpha + \tau_{xy} \sin 2\alpha
    \end{myformula}
    注意：此处符号约定与常见的“拉正压负、顺时针为正”匹配，若你教材中 $\sigma_\alpha$ 表达式第一项为减号，只需将 $\alpha$ 的正方向反向即可。

    \item \textbf{切向平衡方程：}
    沿 $t$ 轴（$n$ 轴顺时针 $90^\circ$）投影：
    \begin{equation*}
        \tau_\alpha dA = \sigma_x (dA \sin\alpha) \cos\alpha - \sigma_y (dA \cos\alpha) \sin\alpha - \tau_{xy} (dA \sin\alpha) \sin\alpha + \tau_{yx} (dA \cos\alpha) \cos\alpha
    \end{equation*}
    化简得：
    \begin{myformula}
        \tau_\alpha = -\frac{\sigma_x - \sigma_y}{2} \sin 2\alpha + \tau_{xy} \cos 2\alpha
    \end{myformula}
\end{enumerate}

上述公式表明，斜截面上的正应力 $\sigma_\alpha$ 和切应力 $\tau_\alpha$ 都是 $\alpha$ 的函数，由此可求极值并确定其所在位置。

\subsection{主应力与主平面}

\begin{mydefinition}{主平面与主应力}
    切应力 $\tau_\alpha$ 等于零的平面称为\textbf{主平面}，其上的正应力称为\textbf{主应力}。主应力按代数值大小排列为 $\sigma_1 \ge \sigma_2 \ge \sigma_3$。
\end{mydefinition}

令切应力公式等于零，得主平面方位角 $\alpha_0$ 满足：
\begin{equation}
    \tan 2\alpha_0 = \frac{2\tau_{xy}}{\sigma_x - \sigma_y}
\end{equation}
由上式可求得两个相差 $90^\circ$ 的角 $\alpha_0$。将 $\alpha_0$ 代入正应力公式，可得两个主应力值：
\begin{myformula}
    \sigma_{\max,\min} = \frac{\sigma_x + \sigma_y}{2} \pm \sqrt{\left( \frac{\sigma_x - \sigma_y}{2} \right)^2 + \tau_{xy}^2}
\end{myformula}
在平面应力状态下，第三个主应力为零（若取 $z$ 方向），因此三个主应力为：
\[
\sigma_1 = \sigma_{\max},\quad \sigma_2 = 0,\quad \sigma_3 = \sigma_{\min} \quad (\text{按代数值排列})
\]

\subsection{最大切应力}

对 $\tau_\alpha$ 求导并令其为零，可得最大切应力所在平面的方位 $\alpha_1$：
\begin{equation}
    \tan 2\alpha_1 = -\frac{\sigma_x - \sigma_y}{2\tau_{xy}}
\end{equation}
比较主平面方位角公式可知，$\tan 2\alpha_1 = -1 / \tan 2\alpha_0$，故最大切应力所在平面与主平面成 $45^\circ$ 夹角。其值为：
\begin{myformula}
    \tau_{\max} = \sqrt{\left( \frac{\sigma_x - \sigma_y}{2} \right)^2 + \tau_{xy}^2} = \frac{\sigma_1 - \sigma_3}{2}
\end{myformula}

\subsection{平面应力状态的正应力不变量}

在平面应力状态下，设单元体的 $x$、$y$ 方向正应力为 $\sigma_x$、$\sigma_y$，切应力为 $\tau_{xy}$，则任意角度 $\alpha$ 截面上的正应力为：
\[
\sigma_\alpha = \frac{\sigma_x+\sigma_y}{2} + \frac{\sigma_x-\sigma_y}{2}\cos2\alpha + \tau_{xy}\sin2\alpha
\]
对于与该截面垂直的截面（角度为 $\alpha+90^\circ$），有：
\[
\sigma_{\alpha+90^\circ} = \frac{\sigma_x+\sigma_y}{2} - \frac{\sigma_x-\sigma_y}{2}\cos2\alpha - \tau_{xy}\sin2\alpha
\]
两式相加得：
\begin{myformula}
    \sigma_\alpha + \sigma_{\alpha+90^\circ} = \sigma_x + \sigma_y = \text{常数}
\end{myformula}
该常数称为\textbf{第一应力不变量}，表明任意两个相互垂直截面上的正应力之和与截面方位无关。

\section{应力圆（莫尔圆）法}

\intro{解析法计算繁琐且直观性差。1882年由德国工程师莫尔 (Otto Mohr) 提出的几何图解法，是分析应力状态最强大的直观工具。}

\begin{myTheorem}{莫尔圆的方程}
    对解析法两公式移项并平方相加，消去 $2\alpha$，得到圆方程：
    \[
    \left( \sigma_\alpha - \frac{\sigma_x + \sigma_y}{2} \right)^2 + \tau_\alpha^2 = \left( \frac{\sigma_x - \sigma_y}{2} \right)^2 + \tau_{xy}^2
    \]
    圆心位于 $\left(\dfrac{\sigma_x + \sigma_y}{2}, 0\right)$，半径 $R = \sqrt{\left(\dfrac{\sigma_x - \sigma_y}{2}\right)^2 + \tau_{xy}^2}$。
\end{myTheorem}

\subsection*{莫尔圆的作图步骤}
\begin{enumerate}
    \item 建立 $\sigma$—$\tau$ 直角坐标系，$\sigma$ 轴向右为正，$\tau$ 轴向下为正（或向上，需统一）。
    \item 在 $\sigma$ 轴上标出 $A(\sigma_x, 0)$、$B(\sigma_y, 0)$。
    \item 过 $A$ 作垂线，取 $\tau_{xy}$ 的正负号（按约定）得点 $D(\sigma_x, \tau_{xy})$；过 $B$ 作垂线，取 $-\tau_{xy}$ 得点 $E(\sigma_y, -\tau_{xy})$。
    \item 连接 $DE$，与 $\sigma$ 轴交于圆心 $C$。
    \item 以 $C$ 为圆心，$CD$ 为半径画圆，即为应力圆。
    \item 圆与 $\sigma$ 轴的两个交点即为两个主应力值 $\sigma_1$ 和 $\sigma_3$。
    \item 从 $D$ 点沿圆周旋转 $2\alpha$（逆时针为正）得到斜截面上的应力点 $P$，其坐标即为 $\sigma_\alpha$ 和 $\tau_\alpha$。
\end{enumerate}

\section{强度理论 (Strength Theories)}

\subsection{破坏准则概述}
材料在外力作用下产生失效（塑性屈服或脆性断裂）的本质原因是什么？历史上提出了四大经典强度理论。对于复杂应力状态，需将多轴应力折算成单向拉伸时的等效应力 $\sigma_{ri}$，并与许用应力比较。

\begin{mydefinition}{第一强度理论（最大拉应力理论）}
    断裂是由最大主应力引起的。当 $\sigma_1$ 达到材料单向拉伸时的强度极限时发生断裂。
    \[
    \sigma_{r1} = \sigma_1 \le [\sigma]
    \]
    适用于脆性材料（如铸铁、石料），且 $\sigma_1$ 为拉应力。
\end{mydefinition}

\begin{mydefinition}{第二强度理论（最大拉应变理论）}
    破坏是由最大主伸长应变引起的。当 $\varepsilon_1 = \frac{1}{E}[\sigma_1 - \mu(\sigma_2 + \sigma_3)]$ 达到单向拉伸断裂时的应变值时发生断裂。
    \[
    \sigma_{r2} = \sigma_1 - \mu(\sigma_2 + \sigma_3) \le [\sigma]
    \]
    适用于脆性材料，但如今已较少使用，因与试验数据吻合不佳。
\end{mydefinition}

\begin{mydefinition}{第三强度理论（最大切应力理论）}
    屈服是由最大切应力引起的。当 $\tau_{\max} = \frac{\sigma_1 - \sigma_3}{2}$ 达到单向拉伸屈服时的 $\tau_s = \frac{\sigma_s}{2}$ 时发生屈服。
    \[
    \sigma_{r3} = \sigma_1 - \sigma_3 \le [\sigma]
    \]
    适用于塑性材料（如低碳钢），偏于安全，但未考虑中间主应力 $\sigma_2$ 的影响。
\end{mydefinition}

\begin{mydefinition}{第四强度理论（形状改变比能理论）}
    屈服是由形状改变比能（畸变能）引起的。当畸变能达到单向拉伸屈服时的畸变能时发生屈服。
    \[
    \sigma_{r4} = \sqrt{\frac{1}{2}\left[(\sigma_1-\sigma_2)^2 + (\sigma_2-\sigma_3)^2 + (\sigma_3-\sigma_1)^2\right]} \le [\sigma]
    \]
    适用于塑性材料，较第三理论更符合试验结果，常用于对塑性材料的安全设计。
\end{mydefinition}

\begin{myhypothesis}{强度理论的选用建议}
    \begin{itemize}
        \item \textbf{脆性材料（铸铁、陶瓷）}：首选第一强度理论，若 $\sigma_1$ 为压应力，则可能发生剪断，可考虑莫尔强度理论。
        \item \textbf{塑性材料（低碳钢、铝合金）}：优先采用第四强度理论，若需保守设计，可采用第三强度理论。
        \item \textbf{复杂应力状态}：当主应力中有拉有压时，需综合考虑材料的拉压性能差异，可查阅莫尔强度理论（经验准则）。
    \end{itemize}
\end{myhypothesis}

